% Source PDF pages 13--16; standard pages 1--4.
% Fidelity pass 05: source-calibrated title-page anchors, exact source line
% membership where Tinos differs from embedded Times New Roman, source emphasis,
% service-group list geometry, and page-local paragraph rhythm.

% -----------------------------------------------------------------------------
% Source PDF page 13 / printed page 1
% -----------------------------------------------------------------------------
\vspace*{1.777in}
{\sffamily\bfseries\fontsize{24pt}{26pt}\selectfont
IEEE Standard for\\
Modeling and Simulation (M\&S)\\
High Level Architecture (HLA)---\\
Framework and Rules\par}

\vspace{6.7pt}
\clause{1. Overview}
\vspace{-3.1pt}
\subclause{1.1 Scope}
This document provides an overview of the High Level Architecture (HLA), defines a family of related\\
HLA documents, and defines the principles of HLA in terms of responsibilities that federates (simulations,\\
supporting utilities, or interfaces to live systems) and federations (sets of federates working together) must\\
uphold.

\vspace{4.2pt}
\subclause{1.2 Purpose}
This document describes the general principles defining the HLA and delineates the set of rules that apply to\\
HLA federations and federates. Each rule is then described, and the rationale for its inclusion is provided.

\vspace{6pt}
Many different classes of simulations exist. Each class has changing application characteristics and needs to\\
be flexibly supported to allow for interoperability and reuse across classes and to limit the need to maintain\\
multiple interoperability approaches. The HLA is a common integrated architecture, which has been

\newpage

% -----------------------------------------------------------------------------
% Source PDF page 14 / printed page 2
% -----------------------------------------------------------------------------
developed to provide a flexible approach for addressing these interoperability and reuse needs. The related\\
standards need to be considered as a set of products because changes in one are likely to have an impact on\\
the others.

\vspace{3.2pt}
\subclause{1.3 High Level Architecture (HLA)}
The HLA is a technical architecture developed to facilitate the reuse and interoperation of simulation\\
systems and assets (e.g., federation managers, data collectors, passive viewers, real world ("live") systems,\\
and other utilities). It was developed in response to the clear recognition of the value of modeling and\\
simulation (M\&S) in supporting a wide variety of applications as well as the need to manage M\&S assets as\\
cost effective tools. The selected definition of an architecture as intended in HLA---``major functional\\
elements, interfaces, and design rules, pertaining as feasible to all simulation applications, and providing a\\
common framework within which specific system architectures can be defined''---is one that is commonly\\
accepted and is consistent with the definition of architecture for computer simulations provided by the IEEE.

\subsubclause{1.3.1 Components of the HLA standard}
\vspace{1.1pt}
The HLA provides a general framework within which developers can structure and describe their simulation\\
systems and/or assets and interoperate with other simulation systems and assets. The HLA consists of three\\
main components. The first is the HLA Framework and Rules Specification (i.e., this standard), which\\
provides a set of ten rules, five applying to \textit{federates} (see 3.1) and five applying to \textit{federations} (see 3.1), that\\
together ensure the proper interaction of federates in a federation and define the responsibilities of federates\\
and federations. The second component is the HLA object model template (OMT) Specification (found in\\
IEEE Std 1516.1\texttrademark-2010).\footnote{Information on references can be found in Clause 2.} It provides a specification for describing \textit{object models} (see 3.1) that define the\\
information produced or required by a simulation application and for reconciling definitions among\\
simulations to produce a common data model for mutual interoperation. The third component is a\\
specification of services that allow simulations to connect to one another, exchange data, and coordinate\\
activities during a distributed runtime execution. This component is the HLA federate interface specification\\
(found in IEEE Std 1516.2\texttrademark-2010), which describes the capabilities and principles of operation of an\\
underlying software infrastructure, called the \textit{runtime infrastructure (RTI)}, that will be required for runtime\\
operation.

\subsubclause{1.3.2 Terminology}
Some terminology is necessary to give an overview of the HLA framework. A federate application is an\\
application that supports the HLA interface to a RTI and that is capable of joining a federation execution.\\
Note that, in addition to a simulation, a federate application may represent an interface to a live system or a\\
support utility such as an event logger or performance monitor. A federate application may be coupled with\\
other software applications and a RTI. A joined federate is a member of a federation execution, actualized by\\
a federate application invoking the \hlaService{Join Federation Execution} service. A federation is a named set of\\
federate applications and a common federation object model (FOM) that are used as a whole to achieve\\
some specific objective. A federation execution is the actual operation, over time, of a set of joined federates\\
that are interconnected by a RTI.

\vspace{8.2pt}
\hlaNote{Formal definitions for each of these terms can be found in the definitions clause of all three HLA\\
specifications.\footnote{Notes in text, tables, and figures are given for information only and do not contain requirements needed to implement the standard.}}

\newpage

% -----------------------------------------------------------------------------
% Source PDF page 15 / printed page 3
% -----------------------------------------------------------------------------
\subsubclause{1.3.3 Object models}
\vspace{0.8pt}
There are three types of object models in the HLA framework. First is a simulation object model (SOM).\\
The SOM is a specification of the types of information that an individual federate application could provide\\
to HLA federations as well as the information that an individual joined federate can receive from other\\
joined federates in HLA federation executions. The SOM may be viewed as a federate application's native\\
expression of its capabilities and data interfaces. The SOM facilitates simulation reuse and provides a\\
starting point for combining the simulation with other participating simulations for a common purpose.\\
Second is a federation object model (FOM). The FOM is a specification defining the information exchanged\\
at runtime to achieve a given set of federation objectives. This information includes object classes, object\\
class attributes, interaction classes, interaction parameters, Management Object Model (MOM) contents,\\
and other relevant information. Third is the MOM, which is a group of predefined constructs that provide\\
support for monitoring and controlling a federation execution. FOMs can be developed from scratch by\\
using relevant elements of the SOMs of the federate applications that will make up the federation or by using\\
existing FOMs or reference FOMs and any combination of these approaches. If existing SOMs or FOMs are\\
used, it is important to understand that there may be parts of these SOMs/FOMs that may not be relevant to\\
other federate applications in the federation and thus will not be incorporated in the FOM. For example, a\\
ship simulation may model boiler pressure, but if a federation needs only to know ship location, boiler\\
pressure may optionally be dropped from the FOM. Also, it may be necessary to reconcile differing units\\
and differing representations among simulations. The OMT provides a standard structure and a set of\\
representational conventions for specifying both SOMs and FOMs.

\subsubclause{1.3.4 Service groups}
\vspace{1.8pt}
The federate interface specification defines the functional interface between federates and the RTI. The RTI\\
provides services to joined federates in a way that is analogous to how a distributed operating system\\
provides services to applications. The services are arranged into seven basic groups:

\vspace{4.1pt}
\begin{hlalist}
    \item Federation Management. These services allow for the coordination of federation-wide activities\\
    throughout the life of a federation execution. Such services include federation execution creation\\
    and destruction, federate application joining and resigning, federation synchronization points, and\\
    save and restore operations.
    \item Declaration Management. These services allow joined federates to specify the types of data they will\\
    supply to, or receive from, the federation execution. This process is done via a set of publication and\\
    subscription services along with some related services.
    \item Object Management. These services support the life-cycle activities of the objects and interactions\\
    used by the joined federates of a federation execution. These services provide for registering and\\
    discovering object instances, updating and reflecting the instance attributes associated with these\\
    object instances, deleting or removing object instances as well as sending and receiving interactions\\
    and other related services. (Note: Formal definitions for each of these terms can be found in the\\
    definitions clause of all three HLA specifications.)
    \item Ownership Management. These services are used to establish a specific joined federate's privilege to\\
    provide values for an object instance attribute as well as to facilitate dynamic transfer of this\\
    privilege (ownership) to other joined federates during a federation execution.
    \item Time Management. These services allow joined federates to operate with a logical concept of time\\
    and to jointly maintain a distributed virtual clock. These services support discrete event simulations\\
    and assurance of causal ordering among events.
    \item Data Distribution Management. These services allow joined federates to further specify the\\
    distribution conditions (beyond those provided via Declaration Management services) for the\\
    specific data they send or ask to receive during a federation execution. The RTI uses this information\\
    to route data from producers to consumers in a more tailored manner.
    \item Support Services. This group includes miscellaneous services utilized by joined federates for\\
    performing such actions as name-to-handle and handle-to-name transformations, the setting of\\
    advisory switches, region manipulations, and RTI startup and shutdown.
\end{hlalist}

\newpage

% -----------------------------------------------------------------------------
% Source PDF page 16 / printed page 4
% -----------------------------------------------------------------------------
The federate interface specification includes programming language-independent definitions of the syntax\\
and semantics of all the services in the interface. It also provides language bindings for Java, C++, and Web\\
Services Definition Language (WSDL).

\subsubclause{1.3.5 Rules}
The HLA rules complete the framework. They require the use of the OMT and Federate Interface\\
specifications as well as sound development principles and practices to ensure the proper operation of a\\
federation execution.

\subsubclause{1.3.6 Implementation independence}
The HLA allows distributed simulation systems to be developed and executed for a range of purposes in a\\
standardized manner. The HLA does not prescribe a specific implementation, nor does it mandate the use of\\
any particular programming language or software. Over time, as technology advances become available,\\
new and different implementations will be possible within the framework of the HLA.

\vspace{3.4pt}
\subclause{1.4 Relationship of HLA and object-oriented (OO) concepts}
\vspace{1.3pt}
Although the HLA OMT is the standardized documentation structure for HLA object models, FOMs and\\
SOMs do not completely correspond to common definitions of object models in OO analysis and design\\
(OOAD) techniques. In the OOAD literature, an object model is described as an abstraction of a system\\
developed for the purpose of fully understanding the system. To achieve this understanding, most OO\\
techniques recommend defining several views of the system. For HLA object models, the intended scope of\\
the system description is much narrower and focuses specifically on requirements and capabilities for\\
federate information exchange. For SOMs, the intent is to describe the public interface of the federate in\\
terms of an identified set of supported HLA object classes and interaction classes. A more complete\\
description of how a federate is designed and functions internally (for example, a traditional OO object\\
model) should be provided via documentation resources other than the SOM. For FOMs, the intent is to\\
describe information exchange that happens during a federation execution.

\vspace{6pt}
Differences between HLA and OOAD principles and concepts also appear at the individual object level. In\\
the OOAD literature, objects are defined as software encapsulations of data and operations (methods). In the\\
HLA, objects are defined entirely by the identifying characteristics (attributes), values of which are\\
exchanged between federates during a federation execution. Any OO-related behaviors and operations that\\
affect the values of HLA object attributes are kept resident in the federates. Although HLA class structures\\
are driven by subscription requirements and FOM growth concerns, class structures in OO systems are\\
typically driven by efficiency and maintainability concerns.

\vspace{6pt}
As discussed above, HLA object classes are described by the attributes that are defined for them. These\\
attributes are, in OO parlance, data members of the class. The attributes, which are abstract properties of\\
HLA object classes, are referred to as \textit{class attributes}. HLA object instances are spawned via an HLA\\
service using an HLA object class as a template. Each attribute contained by an HLA object instance is\\
called an \textit{instance attribute}.

\vspace{6pt}
OO objects interact via message passing, in which one OO object invokes an operation provided by another\\
OO object. HLA objects do not directly interact. It is the federates that interact, via HLA services, by\\
updating instance attribute values or sending interactions. Also, responsibility for updating the instance\\
attributes associated with an HLA object instance can be distributed among different federates in a\\
federation (effectively distributing responsibility for maintaining the HLA object instance's state across the\\
federation), whereas OO objects encapsulate state locally and associate update responsibilities with\\
operations that are closely tied to the object's implementation in an OO programming language.

\newpage

In addition to the stated semantic variations in shared terminology, other differences may also exist. Precise definitions of all HLA terms can be found in Clause 3.
